% SOURCE: Chinese PDF viewer p. 151 (lower part)--p. 153 (end of Chapter 6).
\section{Simple Identical-Particle Systems}
\subsection*{1) Helium atom}
For an atom with $Z$ electrons, let its Hamiltonian be
\[\op H=\op H^{(0)}+\op H^{(1)},\]
\[\op H^{(0)}=\sum_{i=1}^{Z}\left(
-\frac{\hbar^2}{2m}\nabla_i^2-\frac{Ze^2}{r_i}\right),
\qquad
\op H^{(1)}=\frac12\sum_{\substack{j,k=1\\j\ne k}}^{Z}
\frac{e^2}{|\vect r_j-\vect r_k|}.\]
The term $\op H^{(0)}$ includes the kinetic energies of the electrons and their interactions with the nucleus. $\op H^{(1)}$ is the interaction between electrons; the factor $1/2$ removes double counting. Interactions among three or more electrons have been neglected.

Because $\op{\vect L}\cdot\op{\vect S}$ coupling is not considered, the atomic state is
\[\psi(\vect r_1,\ldots,\vect r_Z)
\chi(s_{1z},\ldots,s_{Zz}).\]
For a helium atom, $Z=2$. Since coupling between $\op{\vect s}_1$ and $\op{\vect s}_2$ is not considered, symmetric and antisymmetric spin states were given above. If the interaction between the two electrons is neglected, the spatial wavefunction is
\[\psi(\vect r_1,\vect r_2)
=\psi_{nlm}(\vect r_1)\psi_{n'l'm'}(\vect r_2),
\qquad E=\varepsilon_n+\varepsilon_{n'},\]
where
\[\varepsilon_n=-\frac{m(Ze^2)^2}{2\hbar^2}\frac1{n^2}.\]
For $\op H^{(0)}$, because of identical-particle symmetry, the fully antisymmetric electronic wavefunction is
\[\psi(\vect r_1,\vect r_2,s_{1z},s_{2z})
=\psi_\pm(\vect r_1,\vect r_2)\chi_\mp(s_{1z},s_{2z}),\]
with
\[\psi_\pm(\vect r_1,\vect r_2)=\frac1{\sqrt2}
\left[\psi_{nlm}(\vect r_1)\psi_{n'l'm'}(\vect r_2)
\pm\psi_{nlm}(\vect r_2)\psi_{n'l'm'}(\vect r_1)\right].\]
For the ground state, only the symmetric spatial wavefunction
\[\psi_+(\vect r_1,\vect r_2)
=\psi_{100}(\vect r_1)\psi_{100}(\vect r_2)\]
is possible; there is no antisymmetric spatial wavefunction.

Consider the energy correction induced by $\op H^{(1)}$. In Chapter 4, the ground-state energy and wavefunction were calculated by the perturbation and variational methods. For an excited state, the energy level $E$ is doubly degenerate with states
\[\psi_{nlm}(\vect r_1)\psi_{n'l'm'}(\vect r_2)
\quad\text{and}\quad
\psi_{nlm}(\vect r_2)\psi_{n'l'm'}(\vect r_1).\]
However, its correct zeroth-order wavefunction is $\psi_\pm(\vect r_1,\vect r_2)$, and the first-order correction to the energy is
\[E_\pm=\int\psi_\pm^*(\vect r_1,\vect r_2)
\frac{e^2}{|\vect r_1-\vect r_2|}
\psi_\pm(\vect r_1,\vect r_2)
\,\dd^3r_1\dd^3r_2=I\pm J,\]
where
\begin{align*}
I&=\int|\psi_{nlm}(\vect r_1)|^2
|\psi_{n'l'm'}(\vect r_2)|^2
\frac{e^2}{|\vect r_1-\vect r_2|}\,
\dd^3r_1\dd^3r_2,\\
J&=\int\psi_{nlm}(\vect r_1)\psi_{nlm}^*(\vect r_2)
\psi_{n'l'm'}(\vect r_2)\psi_{n'l'm'}^*(\vect r_1)
\frac{e^2}{|\vect r_1-\vect r_2|}\,
\dd^3r_1\dd^3r_2.
\end{align*}
Here $I$ is the \emph{direct integral}, and $J$ is the \emph{exchange integral}.

For the spatially symmetric wavefunction $\psi_+(\vect r_1,\vect r_2)$, two identical particles are more likely to approach each other and the repulsion is greater. For the spatially antisymmetric wavefunction $\psi_-(\vect r_1,\vect r_2)$, the two identical particles tend to remain farther apart and the repulsion is smaller. Hence $E_+>E_-$.

\subsection*{2) Angular momentum of the atom}
Consider an excited state. The behavior of the system is determined by the two valence electrons, so it can be regarded as a system of two identical valence electrons.

For an atom with two valence electrons, the spin angular momenta and their quantum numbers are
\[\op{\vect s}_1,\qquad\op{\vect s}_2,\qquad
\op{\vect S}=\op{\vect s}_1+\op{\vect s}_2,\]
\[s_1=\frac12,\qquad s_2=\frac12,\qquad S=0,1.\]
The orbital angular momenta and their quantum numbers are
\[\op{\vect L}_1,\qquad\op{\vect L}_2,\qquad
\op{\vect L}=\op{\vect L}_1+\op{\vect L}_2,\]
\[l_1=l,\qquad l_2=l,\qquad
L=|l_1-l_2|,\ldots,l_1+l_2=0,\ldots,2l.\]
The total angular momentum is
\[\op{\vect J}=\op{\vect L}+\op{\vect S},\]
with quantum numbers
\[
J=|L-S|,\ldots,L+S
=\begin{cases}
L,&S=0\quad\text{(antisymmetric spin state)},\\
L-1,L,L+1,&S=1\quad\text{(symmetric spin state)}.
\end{cases}
\]
\textbf{Question.} Are there restrictions on the possible values of $L$ arising from exchange symmetry?

The total wavefunction must be antisymmetric. The spatial wavefunction can be separated into center-of-mass and relative parts $\Phi(\vect R)\Psi(\vect r)$. The center-of-mass factor $\Phi(\vect R)$, where $\vect R=(\vect r_1+\vect r_2)/2$, is symmetric under exchange of the two electrons because $\vect R\to\vect R$. The relative factor is
\[\Psi(\vect r)=R_{NL}(r)Y_{LM}(\theta,\varphi),
\qquad\vect r=\vect r_1-\vect r_2.\]
Exchanging the two electrons sends $\vect r\to-\vect r$, and the exchange operator is equivalent to the parity operator. Thus
\[\op P_{12}Y_{LM}(\theta,\varphi)
=(-1)^L Y_{LM}(\theta,\varphi).\]
For even $L$, the spatial wavefunction is symmetric; for odd $L$, it is antisymmetric. The antisymmetry requirement on the total wavefunction therefore implies
\begin{itemize}
\item symmetric spatial and antisymmetric spin states, with $L$ even and $S=0$;
\item antisymmetric spatial and symmetric spin states, with $L$ odd and $S=1$.
\end{itemize}
In either case $L+S$ is even. Thus exchange symmetry restricts the wavefunction: $L$ and $S$ cannot be combined arbitrarily.

\subsection*{3) Pressure resulting from the Pauli effect}
Electrons in a solid can be approximated as a free Fermi gas.

The phase-space distribution of classical particles is the Boltzmann distribution:
\[f(k)=e^{-(E-\mu)/T},\qquad E=\frac{k^2}{2m}.\]
The phase-space distribution of quantum fermions, resulting from the Pauli exclusion principle, is the Fermi--Dirac distribution:
\[f(k)=\frac1{e^{(E-\mu)/T}+1}.\]
At $T=0$,
\[f(k)=\theta(\mu-E),\qquad
k<k_F=\sqrt{2m\mu},\]
where $k_F$ is the Fermi momentum.

The electron number density is
\[n=2\int\frac{\dd^3k}{(2\pi)^3}f(k)
=\frac1{\pi^2}\int_0^{k_F}k^2\dd k
=\frac{k_F^3}{3\pi^2}.\]
The factor $2$ arises because an electron has two spin states.

The energy density is
\[\varepsilon=2\int\frac{\dd^3k}{(2\pi)^3}E f(k)
=\frac{k_F^5}{10m\pi^2}.\]
From the thermodynamic relation
\[\varepsilon=-P+\mu n,\]
the pressure is
\[P=\frac23\varepsilon.\]

\textbf{Note.} For a classical system the pressure at $T=0$ is $P=0$. The pressure of the free quantum system is due neither to interaction between electrons nor to thermal motion: it arises entirely from the identical-particle effect, namely the Pauli exclusion principle. It is a quantum effect.